\documentclass[11pt]{article}

\usepackage[margin=1in]{geometry}
\usepackage{amsmath,amssymb,amsthm}
\usepackage{mathtools}
\usepackage{bm}
\usepackage{booktabs}
\usepackage{enumitem}
\usepackage{xcolor}
\usepackage{hyperref}
\usepackage{titlesec}
\usepackage{microtype}
\hypersetup{colorlinks=true,linkcolor=blue!50!black,citecolor=blue!50!black,urlcolor=blue!40!black}

\titleformat{\section}{\large\bfseries}{\thesection}{0.6em}{}
\titleformat{\subsection}{\normalsize\bfseries}{\thesubsection}{0.6em}{}

\newcommand{\F}{\mathbf{F}}
\newcommand{\WO}{\mathbf{W}_O}
\newcommand{\WV}{\mathbf{W}_V}
\newcommand{\vv}{\mathbf{v}}
\newcommand{\kk}{\mathbf{k}}
\newcommand{\hh}{\mathbf{h}}
\newcommand{\zz}{\mathbf{z}}
\newcommand{\pp}{\mathbf{p}}
\newcommand{\norm}[1]{\left\lVert #1 \right\rVert}
\newcommand{\TopK}{\mathrm{TopK}}
\newcommand{\softmax}{\mathrm{softmax}}

\title{\textbf{Residual Accounting, Representational Momentum, and \\ White-Box Thought-Anchor Detection in Reasoning Models}\\[0.4em]
\large A Research Proposal}
\author{}
\date{\today}

\begin{document}
\maketitle

\begin{abstract}
\noindent
Reasoning language models produce long chains of thought (CoT) in which a small number of
sentences---\emph{thought anchors}---disproportionately steer the remaining computation
\cite{bogdan2025thought}. Anchors are presently identified \emph{retrospectively} and
expensively, via black-box counterfactual resampling. This proposal asks whether anchors
can be identified \emph{white-box} and, ultimately, \emph{online}, from quantities computed
during a single forward pass. We build on an established per-source-token decomposition of
the attention update \cite{kobayashi2020attention,kobayashi2021incorporating,ferrando2024information}
to define a \emph{delivered residual contribution} graph, propagate influence through it with a
Katz-style recursion, and pair this with a \emph{representational momentum} signal derived from
the trajectory of hidden states. We are explicit throughout about which components are
established prior art, which are incremental, and which are the genuinely open empirical
questions. The central bet of the proposal is a single falsifiable claim: that a cheap,
forward-time signal predicts black-box anchor importance better than the existing attention
baseline ($\rho \approx 0.22$ in \cite{bogdan2025thought}). All experiments are scoped to a
single 1.5B-parameter reasoning model on commodity (RTX TITAN-class) hardware.
\end{abstract}

\section{Background and prior work}

\paragraph{The residual stream and additive decomposition.}
Following the residual-stream view of transformers \cite{elhage2021mathematical}, every component
writes additively into a shared representation, and attention heads act independently, each
adding a result back into the stream. Because attention is linear in the value vectors once the
attention weights are fixed, the attention update can be decomposed by \emph{source token}.

\paragraph{Per-source contributions: attention is not only a weight.}
Kobayashi et al.\ \cite{kobayashi2020attention} observed that attention weights alone do not
describe information flow, because the attention block computes a weighted sum of
\emph{transformed} value vectors; a token with high attention weight but small transformed-value
norm contributes little. Their reformulation, extended to multi-head and to the residual and
normalization layers in \cite{kobayashi2021incorporating}, expresses the layer output as a sum of
per-source contribution functions $\mathbf{F}_i(\mathbf{x}_j)$ with the output projection folded in.
ALTI \cite{ferrando2022measuring}, GlobEnc \cite{modarressi2022globenc}, and DecompX
\cite{modarressi2023decompx} build token-attribution graphs on top of this decomposition and
aggregate contributions across layers (DecompX by propagating decomposed vectors rather than a
rollout recursion). Information Flow Routes \cite{ferrando2024information} carries the construction
to decoder-only LLMs, folding the combined value--output matrix $\mathbf{W}_{OV}=\WV\WO$ and a
linearized LayerNorm into per-(source, head) contributions.

\paragraph{Propagation across the graph: rollout and flow.}
Abnar and Zuidema \cite{abnar2020quantifying} introduced attention rollout (iterated matrix
multiplication of per-layer attention) and attention flow (a max-flow formulation) to approximate
multi-hop, indirect token influence. ALTI and GlobEnc replace raw attention with Kobayashi-style
contributions inside the same rollout machinery. Recent work extends multi-hop, recursive
attribution specifically to \emph{reasoning} traces at token- and sentence-level granularity
\cite{ferrando2024information}. The recursion we use below is, mathematically, a discounted
reachability / Katz--Bonacich centrality \cite{katz1953new} on the contribution DAG.

\paragraph{Thought anchors.}
Bogdan et al.\ \cite{bogdan2025thought} define thought anchors and provide three attribution
methods: (i) a \emph{black-box} method measuring each sentence's counterfactual importance over
100 resampled rollouts; (ii) a \emph{white-box} method aggregating attention between sentence
pairs, identifying ``broadcasting'' sentences attended to by ``receiver'' heads; and (iii) a
\emph{causal} method that suppresses attention toward a sentence and measures the downstream
effect. Crucially, their white-box receiver-head signal correlates with the black-box ground truth
only weakly (Spearman $\rho \approx 0.22$). This number is the baseline this proposal must beat.

\paragraph{Hidden-state dynamics and ``physics'' framings.}
A growing literature treats the sequence of hidden states as a trajectory or dynamical system:
drift-diffusion / regime-switching models of sentence-level trajectories
\cite{statphys2025}, trajectory-extrapolation and momentum-of-language analyses
\cite{trajdyn2026}, and decompositions of reasoning motion into displacement (progress) and
curvature (stability) \cite{semstep2026}. Test-time steering of ``cognitive'' attention heads has
been demonstrated on the same 1.5B model we target \cite{steering2025}.

\paragraph{Causal interventions.}
Activation patching / causal tracing \cite{meng2022locating,wang2023interpretability} and
attention knockout \cite{rai2024how} are standard tools for testing whether a component is
necessary or sufficient for a behavior. Our intervention experiments (Section~\ref{sec:explore})
are instances of these, applied at the granularity of anchors.

\paragraph{Speculative decoding.}
Speculative decoding \cite{leviathan2023fast} uses a small draft model to propose tokens that a
large target model verifies in parallel, preserving the target distribution and yielding
$2$--$3\times$ (up to $\sim$$6\times$ for EAGLE-class methods \cite{li2024eagle}) latency
reductions; it is standard in production serving. Reward-guided variants \cite{liao2025reward}
already switch adaptively between draft and target based on a quality signal. Our
``anchor-gated'' compute-allocation idea (Section~\ref{sec:explore}) is a special case of this
adaptive-switching family, with an interpretability-derived trigger.

\section{Notation and the core object}

Let a reasoning trace have tokens $1,\dots,T$. At layer $\ell$ and head $h$, let
$\alpha^{\ell,h}_{t,s}$ be the attention weight from target token $t$ to source token $s\le t$,
$\vv^{\ell,h}_s$ the value vector, $\kk^{\ell,h}_s$ the key vector, and $\WO^{h,\ell}$ the
head's output projection. The attention update written into token $t$'s residual stream at
layer $\ell$ is
\begin{equation}
\Delta\hh^{\ell,\mathrm{attn}}_t
= \sum_h \WO^{h,\ell}\!\left(\sum_{s\le t}\alpha^{\ell,h}_{t,s}\,\vv^{\ell,h}_s\right)
= \sum_{s\le t}\F^{\ell}_{s\to t},
\label{eq:decomp}
\end{equation}
where the \textbf{delivered residual contribution} from source $s$ to target $t$ is
\begin{equation}
\boxed{\;\F^{\ell}_{s\to t} \;=\; \sum_h \WO^{h,\ell}\!\left(\alpha^{\ell,h}_{t,s}\,\vv^{\ell,h}_s\right).\;}
\label{eq:F}
\end{equation}
This is the vector actually written into $t$ by $s$ at layer $\ell$. It is exactly the per-source
contribution of \cite{kobayashi2021incorporating} (their multi-head $\mathbf F_i(\mathbf x_j)$) and
of \cite{ferrando2024information} (their $\alpha^h f^h(\mathbf x_j)$ with $\mathbf W_{OV}$); we
retain it as a \emph{vector} rather than collapsing to a norm.

\paragraph{LayerNorm caveat.} Equation~\eqref{eq:decomp} is exactly additive for the raw
attention sub-block output. To account for the contribution actually deposited after the
post-attention LayerNorm and residual add, one linearizes LayerNorm as
$\mathrm{LN}(\sum_j \mathbf u_j)=\sum_j L(\mathbf u_j)+\bm\beta$, as in
\cite{kobayashi2021incorporating,ferrando2022measuring}. We adopt that linearization; $\F$ then
denotes the post-LN delivered vector and additivity holds up to the bias term $\bm\beta$.

\section{Method}

\subsection{Edge weights and the delivered-force graph (Construct 1)}

We define a directed, strictly causal ($s<t$; self-contributions excluded) token graph with
scalar edge weights aggregated over layers,
\begin{equation}
E_{s\to t} \;=\; \sum_{\ell=1}^{L}\norm{\F^{\ell}_{s\to t}}_2 ,\qquad s<t .
\label{eq:edge}
\end{equation}
Sentence-level edges aggregate token edges over sentence spans $S_i,S_j$:
\begin{equation}
E_{i\to j} \;=\; \sum_{s\in S_i}\sum_{t\in S_j} E_{s\to t}.
\label{eq:sentedge}
\end{equation}

\paragraph{Top-$K$ pre-pruning with an $\F$-budget.}
Materializing all $\binom{T}{2}$ edges is $O(T^2 L H d)$. We instead keep, for each target $t$,
only its $K$ strongest incoming sources. To avoid the failure mode of pruning by raw attention
(which discards high-norm/low-weight tokens---precisely the case that motivates using $\F$ at
all \cite{guo2024attention}), we prune by a cheap upper bound on the contribution norm,
\begin{equation}
\widehat{E}_{s\to t} \;=\; \sum_{\ell,h}\alpha^{\ell,h}_{t,s}\,\norm{\WO^{h,\ell}\vv^{\ell,h}_s}_2
\;\ge\; \text{(per-head terms of) } E_{s\to t},
\label{eq:budget}
\end{equation}
where $\norm{\WO^{h,\ell}\vv^{\ell,h}_s}$ is precomputed once per source ($O(T L H d)$ total).
We keep $\mathcal{N}(t)=\TopK_K(\{\widehat E_{s\to t}\}_{s<t})$ and compute the exact
Eq.~\eqref{eq:F} norm only on survivors. This is the ``$\F$-budget'': a fixed per-target edge
budget $K$ chosen large enough not to sever genuine multi-hop relay paths.

\subsection{Propagated influence via a Katz-style DP (Construct 1, cont.)}

Direct edges miss \emph{indirect} influence ($s\to m\to t$). We propagate with a discounted
recursion computed right-to-left over the DAG:
\begin{equation}
P_s \;=\; \sum_{t>s} E_{s\to t}\,(1+\lambda P_t),\qquad \lambda\in[0,1),
\label{eq:katz}
\end{equation}
and analogously $P_i=\sum_{j>i}E_{i\to j}(1+\lambda P_j)$ at the sentence level. Eq.~\eqref{eq:katz}
is Katz--Bonacich centrality \cite{katz1953new} on the delivered-force graph; $\lambda$ tunes the
weight of multi-hop paths ($\lambda=0$ recovers direct out-strength). On the pruned graph the DP
costs $O(TK)$---negligible relative to edge construction. \textbf{Top-$K$ anchor selection is
performed on the propagated scores $P_s$, not on direct edges}, because indirect influence is a
global property and is destroyed by pre-DP edge pruning.

\subsection{Representational momentum (Construct 2)}

Let $\zz_t$ be the final-layer hidden state at position $t$. Define representational
\emph{velocity}
\begin{equation}
\vv_t \;=\; \zz_t - \zz_{t-1},
\label{eq:velocity}
\end{equation}
a per-token \emph{stored mass} combining key-norm (addressability) and value--output norm
(write strength),
\begin{equation}
M_s \;=\; \sum_{\ell,h}\norm{\kk^{\ell,h}_s}_2^2\,\norm{\WO^{h,\ell}\vv^{\ell,h}_s}_2^2 ,
\label{eq:mass}
\end{equation}
and a \emph{kinematic momentum} $\pp_t = M_t\,\vv_t$. We distinguish this descriptive momentum
from the \emph{delivered momentum} $\F_{s\to t}$ of Eq.~\eqref{eq:F}, which is causal (it is the
actual write attributable to $s$). The proposed anchor signatures are:
\begin{itemize}[leftmargin=1.4em,itemsep=2pt]
\item \textbf{Reinforcement / stabilizing anchor:} sustained alignment of momentum,
$\cos(\pp_t,\pp_{t+1})\to 1$ over a window, indicating a token whose influence consistently
pushes the trajectory in a fixed direction.
\item \textbf{Pivot anchor:} a high-mass \emph{turn}, i.e.\ a large $\norm{\vv_t}$ accompanied by
a sharp drop in $\cos(\vv_{t-1},\vv_t)$ (high curvature), indicating a redirection of the
reasoning trajectory.
\end{itemize}
The curvature quantity here is the same one used to define ``stability'' in
\cite{semstep2026}; the novel element is the \emph{anchor taxonomy} mapped onto reinforcement vs.\
turn and its coupling to the mass term \eqref{eq:mass}. We treat \eqref{eq:mass} as a hypothesis,
not a settled definition (see Risks).

\paragraph{Predictive test.} As a stronger probe, we test whether momentum predicts future
trajectory:
\begin{equation}
\widehat{\zz}_{t+k} \;=\; \zz_t + \beta\sum_{s\le t} P_s\,\widehat{\F}_s,\qquad
\widehat{\F}_s = \frac{1}{\,\lvert\{t\}\rvert}\sum_{t>s}\F_{s\to t},
\label{eq:predict}
\end{equation}
and whether accumulated momentum aligns with a termination direction $\zz_{\mathrm{EOS}}$
(estimated as the mean final-layer state at end-of-thought), $\cos(\pp_t,\zz_{\mathrm{EOS}})$,
as a candidate online stopping signal.

\subsection{Online detection via running top-$K$ frequency (Construct 4)}

During generation, for each new token $t$ we compute $\mathcal N(t)$ (its top-$K$ contributors)
and increment a running counter $c_s \mathrel{+}= 1$ for every $s\in\mathcal N(t)$. Tokens with
high accumulated $c_s$ (optionally weighted by $P_s$) are flagged as candidate anchors online.
$c_s$ is a streaming, in-degree approximation to $P_s$. This is cheap (it reuses the same
attribution machinery incrementally) but is a heuristic; it must be controlled for the
\emph{position/sink confound} (recency and attention-sink tokens accumulating high counts
without causal importance), which is the same confound flagged for receiver heads in
\cite{bogdan2025thought}.

\section{Establishing ground truth}
\label{sec:groundtruth}

Because ``anchor'' is defined by downstream causal influence, every white-box signal above must be
validated against a ground truth. We use two.

\paragraph{(A) Black-box resampling.} Following \cite{bogdan2025thought}, for each sentence $i$ we
resample $R$ continuations ($R=100$ ideal; $R\in\{20,30\}$ for budget, see Section~\ref{sec:setup})
from the point after $i$, and from a semantically perturbed alternative, and measure the shift in
final-answer distribution. Large shift $\Rightarrow$ anchor. This is the gold standard and the
dominant compute cost.

\paragraph{(B) A Lyapunov / divergence score (cheaper proxy).} As an alternative ground truth
amenable to a single TITAN, we quantify trajectory sensitivity at sentence $i$ by a finite-time
divergence (Lyapunov-style) score: perturb the hidden state (or resample one sentence) and measure
the exponential rate at which trajectories separate,
\begin{equation}
\Lambda_i \;=\; \frac{1}{\tau}\,\log\frac{\norm{\zz^{\text{pert}}_{i+\tau}-\zz_{i+\tau}}}
{\norm{\zz^{\text{pert}}_{i}-\zz_{i}}} ,
\label{eq:lyapunov}
\end{equation}
averaged over a few perturbations. High $\Lambda_i$ marks sentences where the reasoning trajectory
is unstable/sensitive---a dynamical-systems analogue of counterfactual importance, consistent with
the regime-switching view of \cite{statphys2025}. We will report correlation of (B) with (A) to
calibrate it before using it as a stand-in.

\section{Experimental setup}
\label{sec:setup}

\paragraph{Model and hardware.} DeepSeek-R1-Distill-Qwen-1.5B (the \cite{bogdan2025thought} small
model; $L=28$, $H=12$, $d=1536$, $d_h=128$), fp16, single NVIDIA RTX TITAN (24\,GB).
The 1.5B model ($\sim$3\,GB) and a 32k-token value cache ($\sim$2.7\,GB) fit comfortably; the
attribution graph adds $\lesssim$$0.2$\,Tflop per trace, $<$1\% of generation cost and
memory-cheap once edges are collapsed to scalars (Eq.~\eqref{eq:edge}).

\paragraph{Data.} Math and logic problems with verifiable answers (e.g.\ GSM8K, MATH subset),
matching the regimes where anchors are most pronounced \cite{bogdan2025thought}.

\paragraph{Pipeline.}
\begin{enumerate}[leftmargin=1.6em,itemsep=2pt]
\item Generate a CoT; hook attention weights, values, keys, $\WO$ outputs, and final-layer states.
\item Compute ground truth: resampling (A) and/or Lyapunov (B) per sentence
(Section~\ref{sec:groundtruth}).
\item Build the delivered-force graph: precompute $\norm{\WO\vv}$, prune by the $\F$-budget
Eq.~\eqref{eq:budget}, materialize exact edges Eq.~\eqref{eq:F}--\eqref{eq:edge}, run the Katz DP
Eq.~\eqref{eq:katz} for $P_s$/$P_i$.
\item Compute momentum signals Eqs.~\eqref{eq:velocity}--\eqref{eq:mass} and the online counter
$c_s$.
\item \textbf{Primary evaluation:} Spearman correlation of each white-box score
($P_i$; momentum reinforcement/turn; $c_s$) with ground truth (A), against the
$\rho\approx0.22$ receiver-head baseline \cite{bogdan2025thought}.
\end{enumerate}

\section{Exploratory interventions}
\label{sec:explore}

These test \emph{sufficiency} and \emph{manipulability} of anchors. All are instances of
attention knockout / activation patching \cite{rai2024how,wang2023interpretability} and the
attention-suppression method of \cite{bogdan2025thought}, applied at anchor granularity.

\paragraph{Q1: Can one ``skip'' from one anchor to the next?}
(\emph{a}) \emph{Speculative anchoring} on past traces: treat detected anchors as the targets a
draft must hit, and test whether generation can leap between them. (\emph{b}) \emph{Forced
anchoring}: mask attention \emph{from} all non-anchor sentences (retain only anchor$\to$future
edges) and test whether a competent answer survives. Formally, replace $\alpha_{t,s}\!\leftarrow\!0$
for $s\notin\mathcal A$ (anchor set) and renormalize; measure answer accuracy vs.\ the unmasked run.

\paragraph{Q2: Are anchors sufficient? (noise-fill test)}
Replace non-anchor sentence activations with noise (or mean-ablate, as in \cite{rai2024how}) while
preserving anchors; if accuracy is retained, anchors are \emph{sufficient}. Conversely, filling
\emph{anchors} with noise should sharply degrade accuracy (necessity).

\paragraph{Q3: Math heads and learned algorithms.}
On arithmetic problems, isolate the heads carrying the largest $\F$-mass into anchors and test
whether a compact, interpretable subroutine (a ``math head'' circuit) can be identified, in the
spirit of circuit-level CoT analyses \cite{rai2024how}.

\paragraph{Q4: Anchor-gated speculative decoding.}
Draft with the cheap model; verify spans with the expensive model; upon detecting an anchor,
fall back to expensive single-token decoding until the anchor passes. Note this overlaps heavily
with existing adaptive speculative decoding \cite{liao2025reward,leviathan2023fast}: the novelty,
if any, is whether an interpretability-derived trigger beats the draft--target disagreement signal
that speculative decoding already exploits for free. \emph{Feasibility caveat:} a genuinely large
target may not fit alongside the draft on a 24\,GB TITAN without quantization.

\section{Novelty assessment and positioning}

We are explicit about what is and is not new. The summary below reflects a literature review of the
residual-accounting, context-mixing, thought-anchor, hidden-state-dynamics, and speculative-decoding
lines.

\begin{table}[h]
\centering
\small
\begin{tabular}{@{}p{4.4cm}p{3.2cm}p{6.4cm}@{}}
\toprule
\textbf{Component} & \textbf{Status} & \textbf{Closest prior art / note} \\
\midrule
$\F_{s\to t}$ decomposition (Eq.~\ref{eq:F}) & Established &
\cite{kobayashi2021incorporating,ferrando2024information}; we retain the vector. \\
Delivered-force graph + multi-hop DP (Eqs.~\ref{eq:edge},\ref{eq:katz}) & Largely preempted &
ALTI/GlobEnc/DecompX/IFR \cite{ferrando2022measuring,modarressi2022globenc,modarressi2023decompx,ferrando2024information}; recursion is Katz \cite{katz1953new}. Open sliver: sequence-time (not layer) propagation as an \emph{anchor} score. \\
White-box anchor detection (Construct 3) & Goal preempted &
\cite{bogdan2025thought} already detects white-box; only a \emph{better} correlation than $\rho{\approx}0.22$ is a contribution. \\
Momentum taxonomy: reinforce vs.\ pivot (Construct 2) & Freshest, unproven &
Primitives (displacement/curvature) exist \cite{semstep2026,trajdyn2026,statphys2025}; taxonomy + mass coupling not seen. Mass def.\ \eqref{eq:mass} is a hypothesis. \\
Online top-$K$ frequency detector (Construct 4) & Plausibly novel &
Online probing exists; anchor-frequency streaming not seen. Must rule out position/sink confound. \\
Anchor-gated speculative decoding (Q4) & Coupling plausibly novel &
Adaptive switching is done \cite{liao2025reward}; trigger source is the only new part. \\
Lyapunov ground-truth proxy (Eq.~\ref{eq:lyapunov}) & Plausibly novel as anchor GT &
Dynamical-systems framings exist \cite{statphys2025}; use as anchor ground truth not seen. \\
\bottomrule
\end{tabular}
\end{table}

\paragraph{The single load-bearing claim.} Constructs 2, 4, and Q4 all depend on a
\emph{forward-looking} signal predicting a retrospectively-defined quantity. The whole program
reduces to one falsifiable question: \emph{does a cheap, causal, forward-time signal (propagated
delivered force, or representational momentum) predict black-box anchor importance materially
better than the $\rho\approx0.22$ attention baseline?} If yes, the online and compute-allocation
applications follow; if the correlation lands near $0.22$, the program has no result regardless of
novelty. We therefore prioritize this measurement before building the intervention and decoding
machinery.

\section{Risks and limitations}
\begin{itemize}[leftmargin=1.4em,itemsep=2pt]
\item \textbf{Predictive bet may fail.} The core correlation may not beat baseline; this is the
primary risk and is measured first.
\item \textbf{Mass definition is ad hoc.} Eq.~\eqref{eq:mass} multiplies addressability and write
strength; the product form and the squaring are hypotheses. Ablate against keeping the two terms
separate. Note key-norm sign conventions differ across the KV-eviction literature
\cite{guo2024attention}.
\item \textbf{Position/sink confound.} Online frequency counts may recover recency and attention
sinks rather than anchors; requires explicit controls.
\item \textbf{Retrospective definition vs.\ online goal.} True anchorhood is defined by the future;
online detection necessarily predicts it from partial evidence.
\item \textbf{Fast-moving field.} The hidden-state-dynamics literature is producing results
monthly; novelty of Constructs 2/4 should be re-verified immediately before submission.
\end{itemize}

\begin{thebibliography}{99}
\small
\bibitem{bogdan2025thought} P.~C. Bogdan, U.~Macar, N.~Nanda, A.~Conmy.
\emph{Thought Anchors: Which LLM Reasoning Steps Matter?} arXiv:2506.19143, 2025.
\bibitem{kobayashi2020attention} G.~Kobayashi, T.~Kuribayashi, S.~Yokoi, K.~Inui.
\emph{Attention is Not Only a Weight: Analyzing Transformers with Vector Norms.} EMNLP 2020.
\bibitem{kobayashi2021incorporating} G.~Kobayashi, T.~Kuribayashi, S.~Yokoi, K.~Inui.
\emph{Incorporating Residual and Normalization Layers into Analysis of Self-Attention.} EMNLP 2021.
\bibitem{elhage2021mathematical} N.~Elhage et al.
\emph{A Mathematical Framework for Transformer Circuits.} Transformer Circuits Thread, 2021.
\bibitem{ferrando2022measuring} J.~Ferrando, G.~I. G\'allego, M.~R. Costa-juss\`a.
\emph{Measuring the Mixing of Contextual Information in the Transformer (ALTI).} EMNLP 2022.
\bibitem{modarressi2022globenc} A.~Modarressi et al.
\emph{GlobEnc: Quantifying Global Token Attribution by Incorporating the Whole Encoder Layer.} NAACL 2022.
\bibitem{modarressi2023decompx} A.~Modarressi et al.
\emph{DecompX: Explaining Transformers Decisions by Propagating Token Decomposition.} ACL 2023.
\bibitem{ferrando2024information} J.~Ferrando, E.~Voita.
\emph{Information Flow Routes: Automatically Interpreting Language Models at Scale.} EMNLP 2024. arXiv:2403.00824.
\bibitem{abnar2020quantifying} S.~Abnar, W.~Zuidema.
\emph{Quantifying Attention Flow in Transformers.} ACL 2020.
\bibitem{katz1953new} L.~Katz.
\emph{A New Status Index Derived from Sociometric Analysis.} Psychometrika 18(1), 1953.
\bibitem{guo2024attention} Z.~Guo et al.
\emph{Attention Score is Not All You Need (Value-Aware Token Pruning).} EMNLP 2024.
\bibitem{statphys2025} \emph{A Statistical Physics of Language Model Reasoning.} arXiv:2506.04374, 2025.
\bibitem{trajdyn2026} \emph{Trajectory Dynamics in Language-Model Hidden States.} arXiv:2606.05346, 2026 (verify).
\bibitem{semstep2026} \emph{Semantic Step Prediction / Progress and Stability Decomposition.} arXiv:2604.18464, 2026 (verify).
\bibitem{steering2025} \emph{Understanding and Steering the Cognitive Behaviors of Reasoning Models at Test-Time.} arXiv:2512.24574, 2025.
\bibitem{meng2022locating} K.~Meng, D.~Bau, A.~Andonian, Y.~Belinkov.
\emph{Locating and Editing Factual Associations in GPT (ROME / causal tracing).} NeurIPS 2022.
\bibitem{wang2023interpretability} K.~Wang, A.~Variengien, A.~Conmy, B.~Shlegeris, J.~Steinhardt.
\emph{Interpretability in the Wild: a Circuit for Indirect Object Identification (path patching).} ICLR 2023.
\bibitem{rai2024how} D.~Rai, Z.~Yao.
\emph{How to think step-by-step: A mechanistic understanding of chain-of-thought reasoning.} arXiv:2402.18312, 2024.
\bibitem{leviathan2023fast} Y.~Leviathan, M.~Kalman, Y.~Matias.
\emph{Fast Inference from Transformers via Speculative Decoding.} ICML 2023.
\bibitem{li2024eagle} Y.~Li et al.
\emph{EAGLE: Speculative Sampling Requires Rethinking Feature Uncertainty.} ICML 2024.
\bibitem{liao2025reward} B.~Liao et al.
\emph{Reward-Guided Speculative Decoding for Efficient LLM Reasoning.} arXiv:2501.19324, 2025.
\end{thebibliography}

\end{document}
